% Target language: French (France), fr-FR.
% Source: content/en-US.tex, checked against content/zh-CN.tex at upstream commit 0d064f6.
% Translation: AI-assisted; needs native-speaker review before publication.
\StoryParagraph{3}{Le poids d’un seul ovule : 3,5 microgrammes.}

\StoryParagraph{3}{Le poids de cinquante millions de dollars en liquide : 2,5 tonnes.}

\StoryParagraph{3}{Au téléphone depuis le Montage Laguna Beach, Jing Tian m’a demandé le second, en donnant le premier en garantie.}

\StoryParagraph{3}{Elle s’y était rendue pour un prélèvement d’ovocytes. Le bébé, qui devait être porté par une mère porteuse, était attendu en 2027. Si les choses allaient assez vite, il pourrait encore naître sous le signe du Cheval. C’était elle qui en avait parlé la première. À la clinique, elle avait dit : sans les cinquante millions, pas de prélèvement.}

\UnprintedPause


\StoryParagraph{0}{Le Montage Laguna Beach se dresse sur une falaise au-dessus du Pacifique. Nous avions réservé l’étage entier pour trente jours, un million de dollars, afin de préserver son intimité. Elle est partie au bout de huit jours.}

\StoryParagraph{0}{Pendant les vingt-deux jours restants, les femmes de chambre ont continué à monter chaque jour, comme d’habitude. C’est mon assistante qui me l’a raconté. Elle les a vues une fois, lorsqu’elle est allée rendre la clé : une Mexicaine poussait un chariot, entrait dans chaque chambre, changeait les serviettes, essuyait les surfaces, bordait les draps, ressortait et refermait la porte derrière elle. Deux heures pour faire tout l’étage.}

\StoryParagraph{0}{Mon assistante lui a demandé pourquoi elle continuait alors que personne ne séjournait là.}

\StoryParagraph{0}{Elle a répondu que, d’après le registre, quelqu’un occupait les lieux.}

\StoryParagraph{0}{Un Versailles qui avait perdu sa reine.}

\StoryParagraph{0}{J’ai compris soudain que Jing Tian aimait cela. Jing Tian aimait qu’un lieu reste vide à cause d’elle. Que rien ne s’y passe.}

\StoryParagraph{3}{Vide avec le Montage Laguna Beach, un Gulfstream G700 attendait gauchement à Van Nuys tandis que son équipage venait prendre son service chaque jour, comme d’habitude.}

\StoryParagraph{3}{Le huitième jour, elle a passé cet appel. Après cela, personne ne l’a revue.}

\NarrativeBreak


\StoryParagraph{0}{En 2007, quelqu’un dépensait déjà des fortunes pour elle. Cette année-là, je venais d’entrer à l’université de Pékin.}

\StoryParagraph{3}{Bien sûr, ce n’était pas moi.}

\StoryParagraph{3}{Une fenêtre QQ est apparue avec son visage comme égérie. C’était sur ce vieux Lenovo IBM, avec une résolution si mauvaise que tout était flou. J’ai enregistré l’image.}

\StoryParagraph{3}{Dans ma chambre de l’université de Pékin, dont le loyer annuel était de mille yuans, je me suis connecté, je l’ai trouvée sur Xiaonei et je lui ai envoyé une demande d’ajout à ses amis. Xiaonei exigeait un message. J’ai écrit puis effacé, effacé puis réécrit, pendant une quarantaine de minutes. À la fin, j’ai envoyé : « Bonjour. »}

\StoryParagraph{3}{Elle ne l’a pas acceptée.}

\StoryParagraph{3}{Peut-être n’avait-elle même jamais utilisé Xiaonei.}

\StoryParagraph{3}{Dans les années qui ont suivi, quoi que je fasse, je revenais toujours à la même question : jusqu’où devrais-je aller pour qu’elle accepte ? J’y ai pensé quand j’ai gagné mon premier milliard. J’y ai pensé quand j’en ai atteint dix.}

\StoryParagraph{3}{Puis j’ai cessé d’y penser. Renren avait fait faillite depuis longtemps.}

\StoryParagraph{3}{Je n’ai jamais envoyé de seconde demande.}

\StoryParagraph{3}{En 2011, il faisait froid à Pékin. J’ai acheté une doudoune pour cent cinquante yuans. Je me souviens du prix parce que j’ai hésité pendant trois jours.}

\StoryParagraph{3}{La même année, elle a joué dans un film à cent cinquante millions de yuans, \textit{The Warring States}. Sun Honglei, Francis Ng, Kim Hee-sun, Kiichi Nakai et Jiang Wu avaient été réunis pour servir de décor, comme le Montage Laguna Beach et le Gulfstream G700 dix-neuf ans plus tard. Ils ignoraient pourquoi ils étaient là. Puisqu’ils avaient fait le déplacement, il fallait bien que quelque chose arrive. Rien n’est arrivé.}

\StoryParagraph{3}{Puis il y a eu \textit{The Great Wall}, doté d’un budget de cent cinquante millions de dollars. Zhang Yimou. Andy Lau. Matt Damon. Puis Jackie Chan. La liste de l’univers Jing Tian ne cessait de s’allonger.}

\StoryParagraph{3}{À cette époque, j’étais déjà aux États-Unis. J’étudiais, je lançais une entreprise, j’étais de plus en plus occupé. Une seule chose ne changeait pas : chaque fois que je changeais d’appareil, quelles que soient les difficultés, j’y transférais cette image de QQ.}

\StoryParagraph{3}{Dix-neuf ans plus tard, mon nom a rejoint la liste.}

\NarrativeBreak

\StoryParagraph{3}{Nous nous sommes rencontrés pour la première fois au Rosewood Hong Kong. J’avais toujours pensé que je serais nerveux le jour où je la verrais.}

\StoryParagraph{3}{Je ne l’étais pas.}

\StoryParagraph{3}{Elle était plus mince que sur les photos, vêtue sans apprêt, adossée à un canapé. Je l’ai regardée et, soudain, je n’ai plus su ce que j’avais attendu pendant tout ce temps.}

\StoryParagraph{3}{Je l’ai emmenée voir la banane de Cattelan. Elle l’a longuement regardée, puis elle m’a demandé : « Combien ça coûte ? »}

\StoryParagraph{3}{Six millions deux cent mille dollars.}

\StoryParagraph{3}{Elle a dit : « Et quand elle pourrira ? »}

\StoryParagraph{3}{Je n’ai rien répondu. Elle a dit : « Laisse tomber. Allons voir un film. »}

\StoryParagraph{3}{J’ai dit : « D’accord. »}

\StoryParagraph{3}{Elle a dit que personne d’autre que nous deux ne devait se trouver dans la salle.}

\StoryParagraph{3}{J’ai dit : « D’accord. »}

\StoryParagraph{3}{Je suis sorti passer un appel. Quand je suis revenu, elle se tenait encore devant la banane.}

\StoryParagraph{3}{Le film était \textit{Zootopie 2}.}

\StoryParagraph{3}{Une salle de projection pour deux. Ce genre de scène est d’ordinaire réservé au visionnage de contrôle du Département central de la propagande.}

\StoryParagraph{3}{J’ai demandé à mon assistante d’acheter tous les sièges du cinéma.}

\StoryParagraph{3}{J’avais lu des articles qui dénonçaient les faux chiffres du box-office : toutes les places vendues, deux personnes seulement dans la salle. Je comprenais à présent. Ce n’était peut-être pas toujours pour le box-office.}

\StoryParagraph{3}{Il n’était pas difficile d’acheter toutes les places. Le problème venait des billets déjà vendus. Mes assistants et mes gardes du corps sont arrivés en avance et ont attendu à l’entrée. Dès qu’une personne munie d’un billet se présentait, ils lui tendaient des espèces.}

\StoryParagraph{3}{Un couple est arrivé avec son enfant. À la porte, l’enfant brandissait déjà un pot de pop-corn.}

\StoryParagraph{3}{Une assistante leur a tendu une liasse de billets. L’enfant ne comprenait pas pourquoi ils ne pouvaient pas entrer. Les parents ont pris l’argent, remercié, n’ont posé aucune question et sont repartis.}

\StoryParagraph{3}{Ce jour-là, vingt-six personnes ont accepté l’argent.}

\StoryParagraph{3}{Après le film, elle a remplacé sa photo de profil WeChat par l’agent Judy Hopps.}

\StoryParagraph{3}{Ce soir-là, je lui ai tripoté les seins. Elle m’a repoussé.}

\StoryParagraph{3}{J’ai cru qu’elle ne voulait pas. J’ai retiré ma main et j’allais dire quelque chose.}

\StoryParagraph{3}{Elle n’a rien dit. Elle s’est redressée, a sorti une petite pochette de son sac et l’a ouverte. À l’intérieur se trouvait une rangée d’objets, comme une minuscule trousse à outils.}

\StoryParagraph{3}{Elle a dit : « Tes ongles me font mal. »}

\StoryParagraph{3}{J’ai regardé. Je les avais toujours coupés n’importe comment avec un coupe-ongles. Les bords étaient pointus. D’ordinaire, je ne le remarquais même pas.}

\StoryParagraph{3}{Elle a dit : « Viens. »}

\StoryParagraph{3}{Elle a posé ma main sur ses genoux et limé mes ongles un à un.}

\StoryParagraph{3}{D’abord le côté rugueux. Quand elle en a eu fini, elle a changé de face, puis encore une fois. De temps à autre, elle levait un doigt vers la lumière, l’examinait et reprenait.}

\StoryParagraph{3}{J’ai demandé : « Combien de temps ça va prendre ? »}

\StoryParagraph{3}{Elle a dit : « Jusqu’à ce qu’ils ne griffent plus personne. »}

\StoryParagraph{3}{Dix doigts. Cela lui a pris près d’une heure.}

\StoryParagraph{3}{Tout en travaillant, elle m’expliquait : ce côté donne la forme, celui-ci polit et, avec le dernier, il ne faut pas frotter dans les deux sens. Il faut aller dans une seule direction.}

\StoryParagraph{3}{J’ai dit : « D’accord. »}

\StoryParagraph{3}{Quand elle a terminé, elle a retourné ma main, l’a examinée et a passé les doigts dessus.}

\StoryParagraph{3}{Elle a dit : « Voilà. »}

\StoryParagraph{3}{« Vas-y. Touche où tu veux. »}

\StoryParagraph{3}{À un moment, elle a dit quelque chose. Elle a dit : « Si un jour je ne suis plus là, personne ne te polira jamais plus les ongles. »}

\StoryParagraph{3}{J’ai souri. Je n’ai rien dit.}

\StoryParagraph{0}{Trois mois après notre rencontre, j’ai fait constituer un dossier. Ses projets au fil des ans, ses partenaires commerciaux, ses participations dans des sociétés, Jing Guoqing, Tian Xin'ai, tous les comptes liés, l’appartement loué à Pékin, celui acheté à Shanghai et d’autres choses encore. Il a fallu une semaine. Une quarantaine de pages en tout. Il y avait aussi des histoires d’hommes et de femmes. Elles me semblaient les moins importantes.}

\StoryParagraph{0}{Plus tard, elle m’a parlé de certaines choses qui s’y trouvaient. Ce qu’elle m’en a dit ne correspondait pas tout à fait à ce qui était écrit. Je n’ai rien dit.}

\StoryParagraph{0}{De certaines autres, elle ne m’a jamais parlé. Je n’ai jamais posé de question.}

\StoryParagraph{0}{J’ai cru tout ce qu’elle m’a dit.}

\StoryParagraph{0}{J’ai gardé ce dossier. Je l’ai toujours.}

\StoryParagraph{3}{Mais je n’en croyais pas un mot.}

\NarrativeBreak

\StoryParagraph{3}{Elle a dit : « Il y a trop de monde à Hong Kong. On va me reconnaître. Partons quelque part. »}

\StoryParagraph{3}{J’ai dit : « D’accord. »}

\StoryParagraph{3}{Elle a désigné Tenerife sur la carte.}

\StoryParagraph{3}{J’ai dit : « D’accord. »}

\StoryParagraph{3}{Je suis sorti passer un appel. Ce n’est qu’à mon retour que j’ai appris que l’endroit se trouvait dans l’Atlantique, au large de la côte nord-ouest de l’Afrique. Le lendemain à midi, notre groupe de trente personnes est parti. Un A330 nous attendait à l’aéroport.}

\StoryParagraph{3}{Dans la chambre de l’avion, elle m’a demandé : « Tu trouves que mon nom sonne bien ? »}

\StoryParagraph{3}{J’ai dit : « Oui. »}

\StoryParagraph{3}{Elle a dit : « C’est le nom de famille de mon père et celui de ma mère mis bout à bout. Tu ne trouves pas que c’est une drôle de coïncidence ? »}

\StoryParagraph{3}{J’ai dit : « Je l’ai retenu dès le premier jour. »}

\StoryParagraph{3}{Elle m’a pris dans ses bras et a dit : « Ne m’appelle plus Jing Tian. Appelle-moi Maman. »}

\StoryParagraph{3}{« D’accord, Maman. »}

\StoryParagraph{3}{Après l’atterrissage, le commandant de bord est venu me voir. « Patron, l’île a d’abord refusé notre atterrissage. Ils n’avaient encore jamais vu de jet privé aussi grand. »}

\StoryParagraph{3}{Les règles étaient les mêmes sur l’île : personne ne devait se trouver autour de nous. Mes assistants s’en sont chargés. Personne n’a demandé pourquoi. Cette fois, je n’ai pas demandé combien ils étaient.}

\StoryParagraph{3}{Face à l’Atlantique, le Ritz-Carlton avait d’immenses baies vitrées. Elle ne pouvait pas dormir s’il restait la moindre lueur. Les assistants ont recouvert toute la chambre de ruban adhésif noir. Dehors, le jour resplendissait ; dedans, il faisait aussi sombre que dans une cellule d’isolement. Le ruban se décollait au soleil et il fallait le remettre chaque jour.}

\StoryParagraph{3}{Il y avait onze assistants en tout. Ils arrivaient tous les matins à six heures et demie et ne repartaient qu’après notre endormissement. Deux jeunes femmes s’occupaient du ruban : l’une faisait les cadres des fenêtres ; l’autre, les voyants de la climatisation et des prises. Une fois leur travail fini, elles vérifiaient tout trois fois. Elles ont acheté trente rouleaux de ruban, puis vingt de plus. Elles ont remarqué que celui qu’elles posaient après trois heures de l’après-midi se décollait plus facilement parce que le verre était alors plus chaud ; elles se sont donc mises à le faire le matin. Personne ne le leur avait appris. Elles l’avaient découvert seules. Elles en discutaient comme d’un savoir-faire.}

\StoryParagraph{3}{Elle connaissait les noms de ces deux jeunes femmes. Pas moi.}

\StoryParagraph{3}{Nous nous sommes arrêtés au bord de la route, dans le parc national du Teide, pour regarder le coucher du soleil. Les assistants ont sorti des roses de la voiture et ouvert du champagne.}

\StoryParagraph{3}{Maman a contemplé la mer de nuages et a dit : « C’est magnifique. Mais moins que le coucher de soleil à Delingha, dans le Qinghai. »}

\StoryParagraph{3}{Je viens du Qinghai.}

\StoryParagraph{3}{Je ne l’ai pas dit.}

\StoryParagraph{3}{La nuit, nous avons regardé les étoiles depuis le sommet. Malgré nos manteaux matelassés, nous avions encore froid.}

\StoryParagraph{3}{Maman a dit : « Fais un vœu. »}

\StoryParagraph{3}{J’ai demandé : « Qu’est-ce que tu as souhaité ? »}

\StoryParagraph{3}{Maman a dit : « Je ne te le dirai pas. »}

\StoryParagraph{3}{Maman a dit : « C’est magnifique. Mais il fait vraiment trop froid. On va aux Maldives ? »}

\StoryParagraph{3}{J’ai dit : « D’accord. »}

\StoryParagraph{3}{Il faisait moins quinze degrés à Tenerife et vingt-sept aux Maldives. Neuf mille cinq cents kilomètres à vol d’oiseau. Onze heures de vol. Mais l’aéroport des Maldives était trop petit ; en chemin, il nous faudrait donc quitter notre gros-porteur pour un avion plus petit.}

\StoryParagraph{3}{Maman a dit : « Les petits avions font trop de bruit. Les grands sont mieux. »}

\StoryParagraph{3}{Dans la chambre de l’avion qui nous emmenait aux Maldives, Maman m’a confié que son vœu était d’avoir un enfant avec moi. Ce serait obligatoirement par gestation pour autrui.}

\StoryParagraph{3}{Je ne lui ai jamais demandé pourquoi. Je pensais qu’elle avait peur d’avoir mal.}

\StoryParagraph{3}{En plein vol, Maman s’est soudain redressée pour chercher son téléphone. Elle s’est connectée à Starlink et a dit : « Mince. »}

\StoryParagraph{3}{J’ai demandé : « Qu’est-ce qu’il y a ? »}

\StoryParagraph{3}{Elle a dit : « Quelqu’un me menace. Il insiste depuis trois jours et je ne lui ai pas répondu. »}

\StoryParagraph{3}{J’ai dit : « Il veut combien ? »}

\StoryParagraph{3}{Elle a dit : « Tu n’as pas à t’en mêler. »}

\StoryParagraph{3}{Elle a tapé pendant longtemps, effaçant puis modifiant, modifiant puis effaçant. Quand elle a eu fini, elle a posé le téléphone, écran contre ses cuisses. Un moment plus tard, elle l’a retourné pour regarder de nouveau.}

\StoryParagraph{3}{Je lui ai demandé de qui il s’agissait.}

\StoryParagraph{0}{Elle a dit : « Quelqu’un qui travaillait avec moi. Aux finances. Il est allé en prison. Il a porté le chapeau. »}

\StoryParagraph{3}{« Maintenant qu’il est sorti, il n’a nulle part où aller et il tient absolument à s’adresser à moi. »}

\StoryParagraph{3}{J’ai dit : « Alors appelle la police et fais-le arrêter. »}

\StoryParagraph{3}{Elle m’a regardé. « Tu ne comprends pas. »}

\StoryParagraph{3}{J’ai dit : « Tu lui as donné de l’argent ? »}

\StoryParagraph{3}{Elle a dit : « Oui. »}

\StoryParagraph{3}{« Combien ? »}

\StoryParagraph{3}{« Quatre-vingt mille. »}

\StoryParagraph{3}{J’ai dit : « Il t’a fallu trois jours pour réfléchir à ça ? »}

\StoryParagraph{3}{Elle n’a rien dit.}

\StoryParagraph{3}{Elle a à peine parlé pendant le reste de la journée. Elle ne s’est endormie qu’après notre atterrissage.}

\StoryParagraph{0}{La valise de Maman est arrivée la première à Soneva Jani. Les bagages de trente personnes ne tenaient pas dans un seul petit avion ; ils ont été répartis entre trois appareils. Maman n’avait qu’une valise. Cette valise a eu un avion pour elle seule. Maman avait dit de ne pas la mettre avec celles des autres.}

\StoryParagraph{0}{Soneva Jani était célèbre pour sa devise « no news, no shoes ». Maman était rassurée.}

\StoryParagraph{3}{Les majordomes de l’hôtel, les gardes du corps et les assistants étaient vingt-six au total, auxquels s’ajoutaient deux pilotes de drone qui nous suivaient à distance. Les assistants se sont remis à tout recouvrir de ruban adhésif : le voyant de la climatisation, les lumières des prises, encore et encore.}

\StoryParagraph{0}{Le soir, nous avons échangé le récit de nos anciennes relations comme on échange des mots de passe, comme une dynastie qui confronte les noms de ses ères. Maman ne tolérait aucune divergence. Mais elle avait un mot interdit : il ne fallait pas parler du sportif.}

\StoryParagraph{0}{Je ne l’ai pas mentionné, même si je savais que le tatouage dans le dos de Maman était pour lui.}

\StoryParagraph{3}{Il était peut-être trop grand pour être effacé.}

\StoryParagraph{0}{Elle a dit : « Nous formons une famille, maintenant. Mes parents ont eu du mal à m’élever. Tu dois leur verser le cadeau nuptial. »}

\StoryParagraph{0}{J’ai dit : « D’accord. »}

\StoryParagraph{0}{Elle m’a donné deux noms, Jing Guoqing et Tian Xin'ai, ainsi que deux numéros de compte.}

\StoryParagraph{0}{Trente millions. En renminbis.}

\StoryParagraph{0}{J’ai fait le virement. J’avais découvert ces deux noms dans les quarante et quelques pages.}

\StoryParagraph{3}{Un peu plus tard, elle m’a tendu son téléphone. Tian Xin'ai avait répondu par un seul mot : « Reçu. »}

\StoryParagraph{0}{Un jour, Maman a dit : « Tu sais à quoi je pense quand je demande ces choses ? »}

\StoryParagraph{0}{J’ai dit : « À quoi ? »}

\StoryParagraph{0}{Elle a dit : « Je me demande si, un jour, tu diras non. »}

\StoryParagraph{0}{J’ai dit : « Je ne le ferai pas. »}

\StoryParagraph{0}{Elle a dit : « Alors ce n’est pas amusant. »}

\StoryParagraph{3}{Je me suis demandé ce qui pouvait bien m’être impossible.}

\StoryParagraph{0}{Quand l’ennui est devenu trop grand, je me suis remis à l’image QQ de 2007, en essayant de la transférer vers un nouvel appareil. Maman s’est penchée vers moi. « Pourquoi gardes-tu en cachette la photo d’une autre femme ? Qui est-ce ? »}

\StoryParagraph{0}{Je suis resté silencieux. Je n’ai rien dit.}

\StoryParagraph{0}{L’image était trop floue ; elle est donc partie faire autre chose.}

\StoryParagraph{0}{Ce soir-là, j’ai supprimé l’image.}

\StoryParagraph{3}{Le lendemain, je l’ai récupérée sur un ancien téléphone.}

\StoryParagraph{0}{Maman faisait elle aussi défiler son téléphone sans arrêt. Je lui ai demandé ce qu’elle regardait.}

\StoryParagraph{0}{Elle a dit : « Rien. Juste toutes ces conneries pour savoir qui sera en tête du générique. »}

\StoryParagraph{0}{Se disputer la tête du générique : dans un film, une seule vedette peut voir son nom apparaître en premier. Même entre le premier rôle masculin et le premier rôle féminin, un seul peut être cité avant l’autre.}

\StoryParagraph{0}{J’ai dit : « C’est ennuyeux. »}

\StoryParagraph{0}{Elle a dit : « Oui. Et le pire, c’est qu’il se forme parfois des couples de tournage. On passe quarante ou cinquante jours sur un plateau, généralement à jouer face au même homme. On couche ensemble la nuit et, le jour, on se dispute la première place au générique. »}

\StoryParagraph{0}{J’ai dit : « C’est inhumain. »}

\StoryParagraph{0}{Elle a dit : « Oui. »}

\StoryParagraph{0}{« Mais une fois le tournage terminé, peu importe combien de fois on a couché ensemble, on fait comme si rien ne s’était passé. »}

\StoryParagraph{3}{Puis elle a continué à lire pendant une demi-heure.}

\StoryParagraph{0}{Une nouvelle année. Le 2 janvier 2026, à Soneva Jani, j’ai demandé Maman en mariage.}

\StoryParagraph{0}{Je l’avais décidé le matin du 1er janvier. La bague était à Hong Kong. Il n’y avait nulle part où en acheter une sur l’île ; une assistante a donc pris un hydravion pour Malé. Elle est revenue cet après-midi-là avec la bague glissée dans une enveloppe de l’hôtel.}

\StoryParagraph{0}{Maman a dit : « Le diamant est trop petit. »}

\StoryParagraph{0}{Maman a dit : « Ce n’est pas grave. Tu m’en offriras un plus gros plus tard. »}

\StoryParagraph{0}{« Au fait, tu peux venir à Pékin ? »}

\StoryParagraph{0}{Je suis resté silencieux.}

\StoryParagraph{0}{Elle a cru que je ne voulais pas venir.}

\StoryParagraph{0}{Maman a dit : « Pourquoi es-tu si gentil avec moi ? »}

\StoryParagraph{0}{J’ai dit : « Parce que je t’aime bien. »}

\StoryParagraph{0}{Elle a demandé : « Tu étais comme ça avec les autres, avant ? »}

\StoryParagraph{0}{J’ai dit : « Non. »}

\StoryParagraph{3}{Elle a souri sans rien dire.}

\StoryParagraph{0}{Un matin, je me suis réveillé tôt. Elle n’était pas dans le lit.}

\StoryParagraph{0}{Je suis sorti et je l’ai trouvée assise sur la terrasse, face à la mer. Toute la mer était noire. Aucune lumière n’était allumée. L’aube n’était pas encore venue.}

\StoryParagraph{0}{Je lui ai demandé pourquoi elle était debout si tôt.}

\StoryParagraph{0}{Elle a dit : « Je compte les jours. »}

\StoryParagraph{0}{J’ai dit : « Tu comptes quoi ? »}

\StoryParagraph{0}{Elle a dit : « Retourne dormir. »}

\StoryParagraph{0}{Je suis resté assis près d’elle un moment, sans savoir quoi dire. Puis je suis retourné me coucher.}

\StoryParagraph{3}{Quand je me suis réveillé de nouveau, elle était déjà maquillée.}

\StoryParagraph{0}{En quatorze jours aux Maldives, personne ne l’a reconnue.}

\StoryParagraph{0}{Le treizième soir, au restaurant, elle a retiré ses lunettes de soleil. Un peu plus tard, elle a également retiré son chapeau et l’a posé sur la table. Elle a passé tout le repas ainsi. À la table voisine se trouvait un couple d’Allemands âgés. Le serveur était sri-lankais. Personne ne lui a accordé un second regard.}

\StoryParagraph{0}{Après le dîner, elle a remis son chapeau.}

\StoryParagraph{0}{Maman a dit : « Rentrons. »}

\StoryParagraph{0}{J’ai dit : « Nous ne devions pas rester plus longtemps ? »}

\StoryParagraph{3}{Elle a dit : « C’est trop calme. »}

\UnprintedPause

\StoryParagraph{0}{De retour à Hong Kong, Maman a rencontré mes parents.}

\StoryParagraph{0}{Ce soir-là, elle a dit que sa famille devait préparer la rencontre entre nos deux familles et le mariage. Je devais envoyer autant que possible. Toujours à Jing Guoqing et Tian Xin'ai.}

\StoryParagraph{0}{Je n’ai pas demandé pourquoi.}

\StoryParagraph{0}{J’ai dit : « Le plafond quotidien de chaque carte est d’un million. »}

\StoryParagraph{0}{Maman a dit : « C’est trop lent. »}

\StoryParagraph{0}{« Utilise cinq cartes. »}

\StoryParagraph{0}{Cinq cartes, deux comptes, en alternant. À la troisième carte, j’ai hésité. J’ai tout de même terminé tous les virements.}

\StoryParagraph{3}{Maman a dit : « Ce n’est pas assez efficace. »}

\StoryParagraph{0}{J’ai dit : « Faisons venir tous nos parents à Hong Kong pour passer le Nouvel An ensemble. »}

\StoryParagraph{0}{Maman a répondu : « Cette année, il est déjà prévu que nous le passions à Hainan. L’année prochaine. »}

\StoryParagraph{0}{Pour la fête du Printemps, les feux d’artifice au-dessus de Victoria Harbour ont attiré un million de personnes. Elle n’était pas parmi elles.}

\StoryParagraph{0}{À trois heures du matin, j’ai raccompagné mes parents chez eux.}

\StoryParagraph{0}{Les rues étaient vides. Aucune autre voiture n’attendait au feu rouge. Je me suis arrêté malgré tout.}

\StoryParagraph{0}{Maman a dit qu’il valait mieux ne pas nous voir pour le moment. Elle allait partir à Los Angeles pour le prélèvement d’ovocytes. Elle avait d’abord des affaires à régler.}

\StoryParagraph{0}{Maman a demandé de nouveau : « Tu peux venir à Pékin ? »}

\StoryParagraph{0}{Je suis resté silencieux.}

\StoryParagraph{3}{Elle n’a jamais dit ce qu’il y avait à Pékin. Je n’ai jamais demandé.}

\StoryParagraph{0}{Maman a dit que nous devions acheter un logement à Hong Kong et le mettre à son nom.}

\StoryParagraph{0}{J’ai dit : « D’accord. »}

\StoryParagraph{3}{Nous ne nous sommes pas vus non plus pour la Saint-Valentin.}

\StoryParagraph{0}{Un jour, au téléphone, elle m’a dit : « Ce matin, on m’a prélevé huit tubes de sang. »}

\StoryParagraph{0}{J’ai dit : « Autant que ça ? »}

\StoryParagraph{0}{Elle a répondu : « Ça allait. »}

\StoryParagraph{0}{J’ai dit : « Ça a dû être éprouvant. »}

\StoryParagraph{0}{Puis nous avons parlé d’autre chose.}

\StoryParagraph{0}{Ce n’est que plus tard que j’ai appris qu’il fallait être à jeun pour ce type de prise de sang. Elle devait avoir lieu le deuxième ou le troisième jour du cycle menstruel ; il ne fallait pas se tromper de date. Elle était probablement partie vers six heures du matin. L’hiver à Pékin. Il faisait encore nuit. Les examens se sont ensuite multipliés. Je le savais parce qu’une nouvelle ligne était apparue dans le relevé mensuel des paiements de mon assistante. Elle avait figuré pour la première fois l’année précédente. Pas grand-chose, quelques milliers de yuans. Après cela, elle revenait chaque mois, parfois deux ou trois fois. Je validais rapidement. Des sommes pareilles ne m’obligeaient pas à regarder le détail. Ce n’est que plus tard que j’ai appris qu’il s’agissait de prises de sang, d’un bilan de six hormones, d’échographies et de l’AMH. Il fallait mesurer sa réserve ovarienne, savoir combien d’ovocytes pourraient être prélevés et s’il valait la peine de continuer.}

\StoryParagraph{0}{Tout avait été fait à Pékin.}

\StoryParagraph{3}{Je ne sais pas dans quel hôpital. Je ne sais pas si elle y est allée seule ou si quelqu’un l’a accompagnée.}

\StoryParagraph{0}{Le 25 février, elle a atterri à Hong Kong.}

\StoryParagraph{0}{Elle a dit qu’elle n’irait pas à Los Angeles tant qu’elle n’aurait pas vu le logement de Hong Kong.}

\StoryParagraph{0}{Je lui ai envoyé la douzaine de logements que nous avions examinés au cours des six mois précédents. Elle les trouvait tous laids. Je n’en avais visité aucun. Mes assistants avaient passé six mois à les parcourir un par un. Ils me les envoyaient ; je les lui transférais. L’un se trouvait dans les Mid-Levels. Sur la vidéo, on voyait Victoria Harbour. Elle avait été tournée au crépuscule. À un moment, la caméra a vacillé et saisi sur le mur l’ombre de la personne qui filmait. J’ai regardé cette ombre en me demandant : qui est cette personne ? Combien a-t-elle marché aujourd’hui ? Combien de logements a-t-elle visités ? Sait-elle à qui l’on montre celui-ci ?}

\StoryParagraph{3}{Le temps manquait. Elle est tout de même montée à bord du G700 à destination de Los Angeles et a rejoint le Montage Laguna Beach le 28 février.}

\NarrativeBreak

\StoryParagraph{0}{Le huitième jour, Maman a appelé. Chez elle, c’était l’après-midi ; chez moi, le milieu de la nuit.}

\StoryParagraph{0}{Elle a dit : « Cinquante millions. De dollars. »}

\StoryParagraph{0}{J’ai dit : « Laisse-moi réfléchir. »}

\StoryParagraph{0}{Elle n’a rien dit. J’entendais l’océan de son côté.}

\StoryParagraph{0}{Au bout d’un moment, elle a dit : « D’accord. »}

\StoryParagraph{0}{Elle a raccroché.}

\StoryParagraph{3}{Depuis que je la connaissais, je ne lui avais jamais dit « laisse-moi réfléchir ».}

\StoryParagraph{0}{J’ai utilisé Claude Code pour interroger les API et passer en revue toutes mes liquidités. La conclusion était que cela n’aurait aucun effet sur moi. J’ai vérifié de nouveau. J’ai saisi la question dans Claude Code. Claude a répondu : « Ne lui donne pas ces cinquante millions de dollars. »}

\StoryParagraph{3}{J’ai demandé : « Elle ne m’aime plus ? »}

\StoryParagraph{3}{Claude a répondu : « L’amour ne m’intéresse pas et je ne le comprends pas. Mais tu ne peux pas lui donner ces cinquante millions de dollars. »}

\StoryParagraph{3}{Je ne lui avais jamais rien refusé. Pas par générosité. Par peur. La peur de la réponse qui viendrait une seconde plus tard.}

\StoryParagraph{0}{Et maintenant, quelqu’un l’avait dit à ma place.}

\StoryParagraph{0}{J’ai fixé cette phrase pendant longtemps. Ce qui me terrifiait, ce n’était pas que Claude soit intelligent. C’était que Claude n’ait pas hésité un seul instant.}

\StoryParagraph{0}{Il ne connaissait aucun dilemme.}

\StoryParagraph{0}{Je me suis dit que ce devait être quelque chose de supérieur. Les humains en étaient incapables.}

\StoryParagraph{3}{Elle ne m’aimait plus ?}

\StoryParagraph{3}{Le lendemain, elle a rappelé.}

\StoryParagraph{3}{Maman a dit : « L’océan est immense. »}

\StoryParagraph{3}{J’ai répondu : « Ah oui ? »}

\StoryParagraph{3}{Elle a dit : « Cet étage est trop grand. Ça résonne. Quelqu’un est venu faire le ménage ce matin. Je me suis cachée dans la salle de bains en attendant son départ. »}

\StoryParagraph{3}{J’ai dit : « Je peux leur demander de ne plus venir. »}

\StoryParagraph{3}{Elle a dit : « Ce n’est pas la peine. »}

\StoryParagraph{3}{Elle a dit : « La clinique a demandé quand le père du bébé arriverait. »}

\StoryParagraph{3}{J’ai dit : « Je peux envoyer quelqu’un. »}

\StoryParagraph{3}{Elle a dit : « Ce n’est pas la peine. »}

\StoryParagraph{3}{Puis elle a demandé : « Tu as décidé ? »}

\StoryParagraph{3}{Je n’ai pas parlé. Ce n’était pas parce que je réfléchissais. Je ne pensais à rien. Au téléphone, il y avait l’océan de son côté et Victoria Harbour du mien. Aucun de nous ne parlait. J’ai compté. Onze secondes.}

\StoryParagraph{3}{Jing Tian a dit : « Très bien. »}

\StoryParagraph{3}{Jing Tian a dit : « Je comprends. »}

\StoryParagraph{3}{Elle a raccroché.}

\StoryParagraph{3}{J’ai repensé d’innombrables fois à ce qu’elle avait dit d’autre, pendant ces deux jours, en dehors des cinquante millions. L’océan est immense. Cet étage est trop grand. Elle s’est cachée dans la salle de bains. C’est tout. Je me souviens de chaque mot. Ce dont je ne parviens pas à me souvenir, c’est de sa voix, du ton qu’elle avait employé en disant : « Je comprends. »}

\StoryParagraph{0}{Je ne m’en souviens pas. J’ai fini par cesser d’essayer.}

\StoryParagraph{0}{J’ai posé le téléphone face contre la table de chevet. Puis j’ai fait quelque chose que je ne comprends toujours pas : j’ai posé une main sur ma poitrine et j’ai compté.}

\StoryParagraph{0}{Mon cœur battait régulièrement.}

\StoryParagraph{0}{J’ai compté de nouveau. Toujours aussi régulièrement. Un peu plus de soixante battements par minute, comme d’habitude. Je suis resté assis, la main pressée contre la poitrine, tel un médecin examinant quelqu’un d’autre.}

\StoryParagraph{3}{J’ai pensé : c’était donc cela.}

\StoryParagraph{0}{Je n’ai rien fait pendant le reste de la journée. Je suis resté dans la chambre jusqu’à la nuit, puis jusqu’au retour du jour, regardant Victoria Harbour passer du bleu au gris, puis au noir. À six heures du matin, les tours financières de l’autre côté du port se sont allumées une à une. Une assistante a frappé une fois pour demander ce qu’elle devait prévoir pour le dîner. J’ai dit que ce n’était pas nécessaire.}

\StoryParagraph{3}{Je n’ai pas pleuré. Je n’ai cessé d’attendre que les larmes viennent.}

\NarrativeBreak

\StoryParagraph{0}{Plus tard, quand mon assistante a réglé les comptes avec moi, une ligne concernait l’indemnisation de l’hôtel à Tenerife. Après notre départ, les jeunes femmes avaient continué à poser le ruban pendant plusieurs jours. Personne ne leur avait dit qu’elles pouvaient arrêter. Selon l’hôtel, le ruban avait été posé trop longtemps, trop souvent. Les traces sur l’encadrement de la fenêtre ne partaient plus. Il fallait remplacer toute la vitre.}

\StoryParagraph{0}{J’ai dit : « Approuvé. »}

\StoryParagraph{3}{Cette vitre donnait sur l’Atlantique.}

\StoryParagraph{0}{Le carburant de l’avion figurant sur la facture coûtait plus cher que l’étage entier du Montage. J’ai posé la question. Mon assistant a expliqué que l’A330 était conçu pour deux cents personnes. Avec seulement trente passagers, il était trop léger et mal équilibré pour le décollage. Il fallait donc, à chaque fois, ajouter du carburant jusqu’à ce qu’il soit assez lourd. Une fois à destination, l’excédent ne pouvait pas rester à bord. Il fallait le larguer.}

\StoryParagraph{0}{J’ai demandé : « À chaque fois ? »}

\StoryParagraph{3}{Il a répondu : « À chaque fois. »}

\StoryParagraph{0}{Il y avait également une boîte renvoyée par la clinique de Los Angeles. Huit jours d’injections de stimulation ovarienne. Elles devaient rester au frais. Aucune n’avait été ouverte.}

\StoryParagraph{3}{Il y avait aussi l’acompte versé à l’agence de gestation pour autrui : deux cent mille dollars, non remboursables. La mère porteuse avait été choisie en janvier. Californienne, trente-quatre ans, deux enfants. Sa photo figurait dans le dossier. Je ne l’avais jamais regardée. Elle était prête depuis janvier.}

\StoryParagraph{3}{Après le départ de Jing Tian, j’ai retrouvé la photo de 2007 et je l’ai regardée une fois. Dix-neuf années s’étaient écoulées. Les pixels de chaque appareil s’amélioraient, tandis que ceux de cette image semblaient se dégrader. Son visage était flou. Je l’ai agrandie et longuement observée, cherchant la moindre différence entre la femme d’alors et celle qu’elle était devenue. Je n’en ai trouvé aucune. La photo n’avait jamais été assez nette pour que je puisse voir.}

\StoryParagraph{0}{Un jour, je me suis coupé les ongles. Quand j’ai eu fini, j’ai passé un doigt dessus. Ils étaient pointus.}

\StoryParagraph{0}{Je me suis souvenu que cet objet se trouvait encore dans le tiroir, celui qu’elle avait laissé.}

\StoryParagraph{0}{Je l’ai sorti et je l’ai passé deux fois sur mes ongles. Ce n’était pas comme il fallait. Elle avait dit d’aller dans une seule direction. J’avais oublié laquelle.}

\StoryParagraph{0}{Je l’ai rangé.}

\StoryParagraph{3}{Mes ongles griffent encore les gens.}

\StoryParagraph{0}{Un matin, un rappel est apparu sur mon téléphone. Mon assistante l’avait programmé l’année précédente en organisant le calendrier : trois mois avant le terme, préparer la chambre.}

\StoryParagraph{3}{Je l’ai balayé.}

\StoryParagraph{0}{Depuis son départ, je ne cesse de penser à une question.}

\StoryParagraph{0}{Si je lui avais donné l’argent ce jour-là, serait-elle restée ?}

\StoryParagraph{0}{J’y ai longuement réfléchi.}

\StoryParagraph{3}{Plus tard, j’ai ressorti le dossier et je l’ai relu. Une quarantaine de pages.}

\StoryParagraph{0}{Puis je me suis souvenu des feux d’artifice au-dessus de Victoria Harbour, le soir de la fête du Printemps, vus depuis le dernier étage.}

\StoryParagraph{0}{Un million de personnes ont déferlé comme la marée. Les feux d’artifice ont pris fin. Un million de personnes se sont retirées.}

\StoryParagraph{0}{Un million de personnes.}

\StoryParagraph{0}{C’est la même chose chaque année.}

\StoryParagraph{0}{Hong Kong à trois heures du matin :}

\StoryParagraph{3}{Rien ne s’est passé.}
